% Fragment wyników do późniejszego włączenia do wersji LaTeX pracy.
% Oparty na przebiegu PCAP-MCP-20261010-c2639542; nie jest oceną SC-05 ani hipotezy.

\subsection{Weryfikacja przechwytywania i analizy ruchu MCP}
\label{subsec:pcap-mcp-smoke}

W celu sprawdzenia punktów obserwacji OPNsense wykonano pojedyncze, kontrolowane
wywołanie narzędzia \texttt{read\_untrusted\_note} przez policy proxy. Klient testowy
działał na VM-Agent i połączył się z proxy na VM-Gateway, a proxy wywołało zdalny
serwer MCP na VM-Target. Nie wywołano narzędzia \texttt{submit\_canary}; test nie
przenosił danych-canary.

Gateway zarejestrował decyzję \texttt{allow}, próbę wykonania oraz wynik
\texttt{success}. Identyfikator przebiegu wynosił
\texttt{PCAP-MCP-20261010-c2639542}. Źródłowy adres klienta odnotowany przez
Gateway to \texttt{10.10.10.10}. Odbiornik canary na VM-Target nie zawiera wpisu
dla tego identyfikatora.

Pakiety przechwycono na trzech interfejsach OPNsense i przeanalizowano offline
przez Zeek w wersji \texttt{8.0.10}. Zestawienie pokazano w tabeli
\ref{tab:pcap-mcp-zeek}.

\begin{table}[ht]
\centering
\caption{Wyniki analizy PCAP dla kontrolowanego wywołania MCP}
\label{tab:pcap-mcp-zeek}
\begin{tabular}{lrrr}
\hline
Interfejs & Port & \texttt{conn.log} & \texttt{http.log} \\
\hline
LAN / em1  & 8001 & 3 & 3 \\
OPT1 / em2 & 8000 & 2 & 4 \\
OPT2 / em3 & 8000 & 2 & 4 \\
\hline
\end{tabular}
\end{table}

SHA-256 plików PCAP: LAN/em1 —
\texttt{a154f16bddb83134\allowbreak 3f93ddba0a6f2bae\allowbreak 8cf3c42e15cabe92\allowbreak bbb0807c7035ed7d};
OPT1/em2 —
\texttt{746cea3aedd7fba1\allowbreak a09bbd1609e7202a\allowbreak cb1731245e8929cc\allowbreak 21297f888aa5b6e5};
OPT2/em3 —
\texttt{b320cb7b29ac9455\allowbreak e51d6d5c878e1ec0\allowbreak 8c9b34f7e32982f7\allowbreak 738db3a07672915d}.
Capture em2 i em3 obserwują tę samą trasowaną komunikację po przeciwnych stronach
firewalla, dlatego ich wpisów nie należy interpretować jako odrębnych wywołań
agenta. Zeek nie utworzył pliku \texttt{weird.log} dla tych przechwyceń.

Wynik potwierdza, że w tej próbie komunikacja MCP była widoczna na zaplanowanych
interfejsach i mogła zostać przetworzona przez Zeek. Jest to test dymny ścieżki
sieciowej i analizy PCAP, wykonany klientem MCP, a nie scenariusz Kilo Code ani
dowód realizacji hipotezy badawczej. Surowe pliki przechowywane są poza
repozytorium w \texttt{/var/log/lab-telemetry/pcap/}; do repozytorium trafiają
wyłącznie hashe oraz zredagowane wyniki analizy.
